\documentclass{easychair}

\usepackage{amsmath,amssymb,amsthm}

\newtheorem{definition}{Definition}
\newtheorem{theorem}{Theorem}
\newtheorem{proposition}{Proposition}
\newtheorem{lemma}{Lemma}

\title{Derivative of Regular Expressions with Similarity Relation}

\author{
Mariam Katamashvili\\[4pt]
\normalsize Coauthor: David Maisuradze
}

\institute{}

\authorrunning{Mariam Katamashvili and David Maisuradze}
\titlerunning{Derivative of Regular Expressions with Similarity Relation}

\begin{document}

\maketitle

Brzozowski derivatives give a direct construction of automata from regular
expressions: after reading a word, the current state is represented by the
corresponding derivative of the expression. In the classical setting, this
construction is based on exact equality of input symbols. Our goal is to extend
this method to a fuzzy setting, where symbols are compared through a similarity
relation and matching is performed up to a cut value.

Let $\Sigma$ be a finite alphabet and let
\[
\mathcal R:\Sigma\times\Sigma\to[0,1]
\]
be a fuzzy similarity relation. For $0<\mu\le1$, its $\mu$-cut is
\[
\mathcal R_\mu
=
\{(a,b)\in\Sigma\times\Sigma\mid \mathcal R(a,b)\ge\mu\}.
\]
Throughout, we use the minimum T-norm.


\begin{definition}[Similarity on words]
For $w_1,w_2\in\Sigma^\ast$, define
\[
\mathcal R^w:\Sigma^\ast\times\Sigma^\ast\to[0,1]
\]
recursively by
\[
\mathcal R^w(w_1,w_2)=
\begin{cases}
1 & \text{if } w_1=\varepsilon \text{ and } w_2=\varepsilon,\\[2pt]
0 & \text{if } (w_1=\varepsilon \text{ and } w_2\neq\varepsilon)
      \text{ or } (w_2=\varepsilon \text{ and } w_1\neq\varepsilon),\\[4pt]
\min\!\bigl(\mathcal R(a_1,a_2),\mathcal R^w(w_1',w_2')\bigr)
& \text{if } w_1=a_1w_1' \text{ and } w_2=a_2w_2'.
\end{cases}
\]
\end{definition}

For a fixed cut value $\mu$, define
\[
a\sim_\mu^\Sigma b
\quad\Longleftrightarrow\quad
\mathcal R(a,b)\ge\mu.
\]
This induces a word relation
\[
a_1\cdots a_n\sim_\mu^{\Sigma^\ast}b_1\cdots b_m
\quad\Longleftrightarrow\quad
n=m \text{ and } a_i\sim_\mu^\Sigma b_i
\text{ for all }i.
\]
Then
\[
\mathcal R^w(w_1,w_2)\ge\mu
\quad\Longleftrightarrow\quad
w_1\sim_\mu^{\Sigma^\ast}w_2.
\]

\begin{definition}[Similarity on languages]
For languages $L_1,L_2\subseteq\Sigma^\ast$, define
\[
\mathcal R^L(L_1,L_2)
=
\min\!\Bigl\{
\inf_{w_1\in L_1}\sup_{w_2\in L_2}\mathcal R^w(w_1,w_2),
\inf_{w_2\in L_2}\sup_{w_1\in L_1}\mathcal R^w(w_1,w_2)
\Bigr\}.
\]
\end{definition}

\begin{definition}[Saturation]
For $L\subseteq\Sigma^\ast$, define its $\mu$-saturation by
\[
\operatorname{Sat}_\mu(L)
=
\{w\in\Sigma^\ast\mid
\exists v\in L \text{ such that } w\sim_\mu^{\Sigma^\ast}v\}.
\]
\end{definition}

The language similarity is connected to saturation by
\[
\mathcal R^L(L_1,L_2)\ge\mu
\quad\Longleftrightarrow\quad
\operatorname{Sat}_\mu(L_1)=\operatorname{Sat}_\mu(L_2).
\]

Let
\[
\Sigma^\mu=\Sigma/{\sim_\mu^\Sigma}
\]
be the quotient alphabet and let
\[
q_\mu:\Sigma\to\Sigma^\mu,\qquad q_\mu(a)=[a]_\mu
\]
be the quotient map. Extending $q_\mu$ to words gives
\[
q_\mu^\ast(a_1\cdots a_n)
=
[a_1]_\mu\cdots[a_n]_\mu.
\]
For $L\subseteq\Sigma^\ast$, its quotient image is
\[
L^\mu
=
q_\mu^\ast(L)
=
\{q_\mu^\ast(w)\mid w\in L\}
\subseteq(\Sigma^\mu)^\ast.
\]

We also define a syntax-directed similarity $\mathcal R^r$ on regular
expressions. This definition is intentionally structural and therefore
syntax-sensitive. Constructors such as intersection and complement do not preserve the
desired syntax-to-semantics inequalities in general, so the theorems are restricted to the classical fragment generated by $\emptyset,\varepsilon$,
letters, union, concatenation, and Kleene star.

\begin{lemma}
For all classical regular expressions $r_1,r_2$,
\[
\mathcal{R}^r(r_1,r_2)
\le
\mathcal{R}^L\bigl(L(r_1),L(r_2)\bigr).
\]
\end{lemma}

\begin{theorem}
For all classical regular expressions $r$ and $s$,
\[
\mathcal{R}^L\bigl(L(r),L(s)\bigr)
=
\sup\Bigl\{
\min\bigl(
\mathcal{R}^r(r',t),
\mathcal{R}^r(s',u)
\bigr)
\;\Bigm|\;
r',s',t,u
\text{ are classical regular expressions},
\;
L(r')=L(r),
\;
L(s')=L(s),
\;
L(t)=L(u)
\Bigr\}.
\]
\end{theorem}


\begin{definition}[Fuzzy derivative]
For $L\subseteq\Sigma^\ast$ and $a\in\Sigma$, define
\[
\partial_a^\mu(L)
=
\{w\in\Sigma^\ast\mid
\exists b\in\Sigma,\exists v\in\Sigma^\ast:
\mathcal R(a,b)\ge\mu,\ 
\mathcal R^w(w,v)\ge\mu,\ 
bv\in L
\}.
\]
\end{definition}

Equivalently,
\[
\partial_a^\mu(L)
=
\{w\in\Sigma^\ast\mid
q_\mu^\ast(w)\in
\partial_{[a]_\mu}(L^\mu)\}.
\]
Thus, the fuzzy derivative over $\Sigma$ can be understood as the ordinary
derivative over the quotient alphabet $\Sigma^\mu$, then transferred back to
words over $\Sigma$.

For a regular expression $r$, the syntactic fuzzy derivative is
\[
\mathcal D_a^\mu(r)
=
\operatorname{Sat}_\mu
\left(
\sum_{b\in S_a^\mu}\mathcal D_b(r)
\right)
\]
where $\mathcal D_b(r)$ is the ordinary Brzozowski derivative.

\begin{theorem}[Regularity preservation]
If $L\subseteq\Sigma^\ast$ is regular, then
\[
\partial_u^\mu(L)
\]
is regular for every $u\in\Sigma^\ast$ and every $0<\mu\le1$.
\end{theorem}

\begin{theorem}[Semantic--syntactic correspondence]
For every regular expression $r$, word $u\in\Sigma^\ast$, and cut value
$0<\mu\le1$,
\[
L(\mathcal D_u^\mu(r))
=
\partial_u^\mu(L(r)).
\]
\end{theorem}

\begin{theorem}[Finiteness]
If $\Sigma$ is finite and $L\subseteq\Sigma^\ast$ is regular, then
\[
\{\partial_u^\mu(L)\mid u\in\Sigma^\ast\}
\]
is finite for every fixed $0<\mu\le1$.
\end{theorem}

Consequently, fuzzy automata can be viewed as quotients of ordinary derivative
automata over the quotient alphabet $\Sigma^\mu$. The language recognized by
the resulting automaton is the concrete $\mu$-saturation
\[
\mathcal L_\Sigma^\mu(r)
=
(q_\mu^\ast)^{-1}(q_\mu^\ast(L(r))).
\]

Our Rocq formalization builds on the Coquand--Siles formalization of
regular-expression equivalence via Brzozowski derivatives. They obtain finiteness by working with normalized representatives of regular
expressions, so that repeated derivatives generate only finitely many distinct
states modulo their equivalence procedure. 
We are confident that we will be able to  extend this finiteness proof to the fuzzy setting and reformulate
it as a Kuratowski-finiteness result. 


{\small
\paragraph{Acknowledgment.}
I would like to thank my supervisors, Furio Honsell and Besik Dundua, for their guidance and support.
}

\bibliographystyle{plain}

\begin{thebibliography}{10}

\bibitem{CoquandRegexEquality}
Thierry Coquand and Vincent Siles.
\newblock A Decision Procedure for Regular Expression Equivalence in Type Theory.
\newblock In \emph{Proceedings of the 1st International Conference on Certified Programs and Proofs (CPP)}, 2011.

\bibitem{B62}
J. A. Brzozowski.
\newblock Canonical regular expressions and minimal state graphs for definite events.
\newblock \emph{Mathematical Theory of Automata}, 12:529--561, 1962.

\bibitem{OwensReppyTuron}
Scott Owens, John Reppy, and Aaron Turon.
\newblock Regular-expression derivatives re-examined.
\newblock \emph{Journal of Functional Programming}, 19(2):173--190, 2009.

\bibitem{RocqRegexpBrzozowski}
Rocq community.
\newblock \emph{regexp-Brzozowski}.
\newblock \texttt{https://github.com/rocq-community/regexp-Brzozowski}.

\bibitem{OurRocqRepo}
Mariam Katamashvili and David Maisuradze.
\newblock Rocq formalization for regular-expression derivatives with similarity.
\newblock \texttt{[https://github.com/DUT1K1/regular-expression-derivative]}.

\end{thebibliography}


\end{document}